% answers_short.tex —— 1994 年·全国卷·文科 · 简短答案速查
% 由 tools/make_short_answers.py 自 answers.tex 的 \gksolhead/\gkanswer 提取，
% 并与 1994/1994_数学真题_全国文_答案要点.md 逐题比对后生成，勿手改（改源 MD 后重跑）。

\begin{gknotice}
\textbf{使用说明}：本节只给最终答案便于快速核对，完整推导见解析版。答案为编写者独立推导，\textbf{非官方标准答案}。
\end{gknotice}

\gkpart{选择题}
\begin{center}\begin{tabular}{cccccccccccccccc}
\toprule
题号 & 1 & 2 & 3 & 4 & 5 & 6 & 7 & 8 & 9 & 10 & 11 & 12 & 13 & 14 & 15 \\
\midrule
答案 & C & D & B & A & B & D & B & A & A & C & C & B & D & D & C \\
\bottomrule
\end{tabular}\end{center}

\gkpart{填空题}
\begin{itemize}[leftmargin=2.6em]
  \item[\textbf{16.}] \(-189\)
  \item[\textbf{17.}] 准线 \(x=3\)；圆 \((x-2)^2+y^2=1\)
  \item[\textbf{18.}] \(-\dfrac34\)
  \item[\textbf{19.}] \(\dfrac{2\sqrt2}{3}\pi\)
  \item[\textbf{20.}] \(\dfrac{a_1+a_2+\cdots+a_n}{n}\)
\end{itemize}

\gkpart{解答题}
\begin{itemize}[leftmargin=2.6em]
  \item[\textbf{21.}] \(-\sqrt2\)
  \item[\textbf{22.}] （条件结论）当 \(a>1\) 时 \(\dfrac12[f(x_1)+f(x_2)]\le f\!\left(\dfrac{x_1+x_2}{2}\right)\)；当 \(0<a<1\) 时不等号反向；两种情形均当且仅当 \(x_1=x_2\) 取等
  \item[\textbf{23.}] (1) 见解析（\(AB_1\parallel\) 平面 \(DBC_1\)，用线面平行判定）；(2) \(\sqrt3\)
  \item[\textbf{24.}] \(\lambda=1\) 时为直线 \(x=\dfrac54\)；\(\lambda\ne1\) 时为圆 \(\left(x-\dfrac{2\lambda^2}{\lambda^2-1}\right)^2+y^2=\dfrac{3\lambda^2+1}{(\lambda^2-1)^2}\)
  \item[\textbf{25.}] 见解析（由 \(S_n-S_{n-1}=a_n\) 推出 \(a_{n+1}+a_{n-1}=2a_n\)）
\end{itemize}
